\documentclass[10pt,a4paper]{article}

\usepackage[utf8]{inputenc}
\usepackage[margin=0.7in]{geometry}
\usepackage{longtable}
\usepackage{array}
\usepackage{booktabs}

\title{\textbf{Extraction Results - New RAI Prompts (All 8 Datasets)}}
\author{Claude Sonnet 4.5}
\date{December 16, 2025}

\begin{document}

\maketitle

\section*{Summary}
\begin{tabular}{|l|c|c|c|}
\hline
\textbf{Dataset} & \textbf{Fill Rate} & \textbf{Fields Filled} & \textbf{Time} \\
\hline
MLS & 63.3\% & 19/30 & 36.94s \\
FLORES & 66.7\% & 20/30 & 47.84s \\
CIFAR & 43.3\% & 13/30 & 23.47s \\
Visual Genome & 70.0\% & 21/30 & 43.43s \\
MSCOCO & 66.7\% & 20/30 & 48.95s \\
MMLU & 66.7\% & 20/30 & 33.54s \\
MMMU & 56.7\% & 17/30 & 38.90s \\
MathVista & 76.7\% & 23/30 & 53.69s \\
\hline
\textbf{Average} & \textbf{63.8\%} & \textbf{19.1/30} & \textbf{40.8s} \\
\hline
\end{tabular}

\newpage

\section*{1. MLS (Multilingual LibriSpeech)}
\textbf{Fill Rate:} 63.3\% (19/30 fields) \quad \textbf{Time:} 36.94s

\begin{longtable}{|p{4cm}|p{12cm}|}
\hline
\textbf{Field} & \textbf{Value} \\
\hline
\endfirsthead
\hline
\textbf{Field} & \textbf{Value} \\
\hline
\endhead
\hline
\endfoot

sc:name & Multilingual LibriSpeech (MLS) \\
\hline
sc:description & A large multilingual corpus suitable for speech research derived from read audiobooks from LibriVox, consisting of 8 languages, including about 44.5K hours of English and a total of about 6K hours for other languages. \\
\hline
sc:url & http://www.openslr.org \\
\hline
sc:license & public domain \\
\hline
sc:creator & Vineel Pratap, Qiantong Xu, Anuroop Sriram, Gabriel Synnaeve, Ronan Collobert \\
\hline
sc:publisher & Facebook AI Research \\
\hline
sc:datePublished & 2020-12-22 \\
\hline
sc:inLanguage & en, de, nl, es, fr, pt, it, pl \\
\hline
cr:citeAs & null \\
\hline
cr:isLiveDataset & null \\
\hline
rai:dataCollection & The dataset is derived from read audiobooks from LibriVox. Audio files were downloaded using LibriLight open source library tools. All audio data was downsampled from 48kHz to 16kHz for further processing. Audio files were segmented into 10-20 second segments using trained acoustic models with Time-Depth Separable (TDS) Convolutions with Auto-Segmentation Criterion (ASG) loss. \\
\hline
rai:dataCollectionType & Web Scraping, Software Collection \\
\hline
rai:dataCollectionMissingData & null \\
\hline
rai:dataCollectionRawData & LibriVox audiobooks, which are mostly based on Project Gutenberg text data. For English, approximately 60K hours of audiobooks were read from four major website domains: gutenberg.org, archive.org, ccel.org and hathitrust.org. For other languages, text sources were extracted from various domains using parsers, manual copying from browsers, or extraction from .pdf/.epub books using pdftotext tool. \\
\hline
rai:dataCollectionTimeframe & null \\
\hline
rai:dataImputationProtocol & null \\
\hline
rai:dataManipulationProtocol & Numbers in the text were replaced by aligning the matched text from the book with the pseudo label and replacing the numbers in the matched text with the corresponding aligned words from the pseudo label. Heuristics were used to deal with hyphens and apostrophes, including removing hyphens used for end-of-line hyphenation and replacing rare words with hyphens or apostrophes based on frequency thresholds. \\
\hline
rai:dataPreprocessingProtocol & Text normalization was performed including NFKC normalization, removal of unwanted characters like punctuations, subscript/superscripts, emojis, escape symbols. Hyphens used for end-of-line hyphenation were removed and word parts were joined. Valid unicode characters were filtered based on language orthography. Audio segmentation was performed to split long audio files into 10-20 second segments based on silence intervals. Pseudo labels were generated using beam-search decoding with a 4-gram language model. Transcript retrieval involved splitting source text into overlapping documents of 1250 words, using TF-IDF similarity score on bigrams, and performing Smith-Waterman alignment. Candidate transcripts with WER >40\% between candidate transcript and pseudo label were filtered out. \\
\hline
rai:dataAnnotationProtocol & Human evaluation of transcripts in development and test sets was performed. Human transcribers listened to audio files while looking at the transcript and were asked to correct only clear and obvious errors to avoid adding potential bias. \\
\hline
rai:dataAnnotationPlatform & null \\
\hline
rai:dataAnnotationAnalysis & WER between generated transcript and human rated transcript was measured. For example, English dev: 4.55\%, test: 4.54\%; German dev: 3.6\%, test: 3.9\%; Dutch dev: 7.24\%, test: 8.61\%; French dev: 3.33\%, test: 4.45\%; Spanish dev: 1.93\%, test: 3.07\%; Italian dev: 6.8\%, test: 5.37\%; Portuguese dev: 12.64\%, test: 12.07\%; Polish dev: 5.67\%, test: 5.12\%. \\
\hline
rai:annotationsPerItem & null \\
\hline
rai:annotatorDemographics & null \\
\hline
rai:machineAnnotationTools & Acoustic models trained using Time-Depth Separable (TDS) Convolutions with Auto-Segmentation Criterion (ASG) loss were used for audio segmentation and pseudo label generation. wav2letter@anywhere framework was used for streaming inference. A gender classifier (SVM with RBF kernel trained on LibriSpeech data) was used to label speaker gender with 95\% test accuracy on LibriSpeech and 96\% and 94\% accuracy on Dutch and Polish respectively. \\
\hline
rai:dataReleaseMaintenancePlan & null \\
\hline
rai:personalSensitiveInformation & Gender information of speakers was collected and used for balancing development and test sets. \\
\hline
rai:dataSocialImpact & null \\
\hline
rai:dataBiases & null \\
\hline
rai:dataLimitations & The dataset is read-speech only, derived from audiobooks. For Polish, the OOV rate is 13\% which negatively impacts WER when performing LM decoding with constrained lexicon. The pseudo label generation process may not always have correct transcription of the audio, which affects the number replacement procedure. \\
\hline
rai:dataUseCases & The dataset is suitable for Automatic Speech Recognition (ASR) research and Text-To-Speech (TTS) research. It includes train, development and test splits. Limited supervision datasets are provided (10 hour set, 1 hour set, and six 10-minute sets) to create a standard benchmark for low resource training. The dataset can be used for training, testing, and validation of ASR models. Training sets are interchangeable and complementary with LibriSpeech for English. \\
\hline
\end{longtable}

\newpage

\section*{2. FLORES (FLORES-101)}
\textbf{Fill Rate:} 66.7\% (20/30 fields) \quad \textbf{Time:} 47.84s

\begin{longtable}{|p{4cm}|p{12cm}|}
\hline
\textbf{Field} & \textbf{Value} \\
\hline
\endfirsthead
\hline
\textbf{Field} & \textbf{Value} \\
\hline
\endhead
\hline
\endfoot

sc:name & FLORES-101 \\
\hline
sc:description & FLORES-101 is an evaluation benchmark consisting of 3001 sentences extracted from English Wikipedia and covering a variety of different topics and domains. These sentences have been translated in 101 languages by professional translators through a carefully controlled process. The resulting dataset enables better assessment of model quality on the long tail of low-resource languages, including the evaluation of many-to-many multilingual translation systems, as all translations are multilingually aligned. \\
\hline
sc:url & https://dl.fbaipublicfiles.com/flores101/dataset/flores101\_dataset.tar.gz \\
\hline
sc:license & null \\
\hline
sc:creator & Naman Goyal, Cynthia Gao, Vishrav Chaudhary, Peng-Jen Chen, Guillaume Wenzek, Da Ju, Sanjana Krishnan, Marc'Aurelio Ranzato, Francisco Guzm\'{a}n, Angela Fan \\
\hline
sc:publisher & Facebook AI Research \\
\hline
sc:datePublished & 2021-06-06 \\
\hline
sc:inLanguage & 101 languages \\
\hline
cr:citeAs & null \\
\hline
cr:isLiveDataset & null \\
\hline
rai:dataCollection & All source sentences were extracted from multiple Wikimedia sources: Wikinews (international news articles), Wikijunior (age-appropriate nonfiction books for children), and WikiVoyage (travel guide with articles about travel tips, food and destinations). The sentence selection process consisted of selecting an article at random from each source, and then manually selecting a few (typically between 3 and 5) contiguous sentences from each article. The selection process was performed by 10 different annotators in the lab, 6 male and 4 female; with different roles in research: researchers (scientists and engineers) and program/project managers; originally coming from different regions of the world: East Asia, South Asia, Southern Europe, Latin America and North America. \\
\hline
rai:dataCollectionType & Manual Human Curator \\
\hline
rai:dataCollectionMissingData & null \\
\hline
rai:dataCollectionRawData & English Wikipedia articles from Wikinews, Wikijunior, and WikiVoyage \\
\hline
rai:dataCollectionTimeframe & null \\
\hline
rai:dataImputationProtocol & null \\
\hline
rai:dataManipulationProtocol & null \\
\hline
rai:dataPreprocessingProtocol & The sentence selection process consisted of selecting an article at random from each source, and then manually selecting a few (typically between 3 and 5) contiguous sentences from each article, avoiding segments with very short or malformed sentences. To avoid bias coming from the document structure, paragraphs were carefully selected from either the beginning, middle or end of the article. The location selection was balanced to be equally distributed across the whole corpus --- roughly one third of paragraphs were sampled from the beginning of the article, one third from the middle, and so on. For each sentence, Wikipedia URL, topic, and boolean flags were extracted to indicate whether the sentence contained entities linked to other Wikipedia pages and images. \\
\hline
rai:dataAnnotationProtocol & For each target language, all source sentences are sent to a translation Language Service Provider (LSP). Once sentences are translated, the data is sent to different translators within the LSP for editing and then moves on to automated quality control steps. If any of the checks fail, the LSP has to re-translate until all verification is passed. Afterwards, translations are sent to an evaluation LSP that performs quality assessment, providing a translation quality score and constructive linguistic feedback both on the sentence and language levels. If the score is below the accepted threshold (90\%), translations together with the assessment report are sent back to the translation LSP for re-translation. The translation quality score is determined based on the number of identified errors by the evaluation LSPs. The following errors are examined: grammar, punctuation, spelling, capitalization, addition or omission of information, mistranslation, unnatural translation, untranslated text, and register. Each error is also associated with a severity level, between minor, major, and critical. \\
\hline
rai:dataAnnotationPlatform & Language Service Providers (LSPs) \\
\hline
rai:dataAnnotationAnalysis & Translation quality was assessed through a Translation Quality Score, calculated per language on a 0 to 100 scale. The translation quality score is determined based on the number of identified errors by the evaluation LSPs. Based on tallying different error types (grammar, punctuation, spelling, capitalization, addition or omission of information, mistranslation, unnatural translation, untranslated text, and register), the overall final score is determined. The acceptable translation quality score was set to 90\%. The largest error category across all languages was mistranslation. Error categories with few errors include register, grammar, and punctuation. \\
\hline
rai:annotationsPerItem & null \\
\hline
rai:annotatorDemographics & The sentence selection process was performed by 10 different annotators in the lab, 6 male and 4 female; with different roles in research: researchers (scientists and engineers) and program/project managers; originally coming from different regions of the world: East Asia, South Asia, Southern Europe, Latin America and North America. \\
\hline
rai:machineAnnotationTools & Automatic checks were implemented to ensure translation quality: (i) language identification, (ii) checking whether the translation is a copy of the source sentence; (iii) checking whether the translation has significantly different length, (iv) checking translation fluency according to a language model, (v) and checking whether the translation is a copy of the translation produced by publicly available translation engines. A heuristic was proposed to detect and reject professional translations that are likely copies from translation engines based on their sentence-level similarity using sentence-level SentencePiece BLEU (spBLEU). \\
\hline
rai:dataReleaseMaintenancePlan & The dev and devtest splits are publicly downloadable. The test set will not be released, but will be available via a publicly available evaluation server. The primary motivation for keeping the test set available only through an evaluation server is to guarantee equivalent assessment of models and reduce overfitting to the test set. \\
\hline
rai:personalSensitiveInformation & null \\
\hline
rai:dataSocialImpact & The majority of the world's population speak low-resource languages and would benefit from improvements in translation quality on their native languages. FLORES-101 enables better assessment of model quality on the long tail of low-resource languages, including the evaluation of many-to-many multilingual translation systems. By publicly releasing such a high-quality and high-coverage dataset, the authors hope to foster progress in the machine translation community and beyond. \\
\hline
rai:dataBiases & null \\
\hline
rai:dataLimitations & While translationese (or overly literal or awkward translations) has known idiosyncrasies, the authors conjecture that these effects are rather marginal when evaluating models in low-resource languages, where current MT systems produce many severe mistakes. All source sentences were extracted from English Wikipedia, meaning that source sentences not in English are produced by human translators. The dataset is suitable for many-to-many evaluation with the caveat that source sentences not in English are produced by human translators. \\
\hline
rai:dataUseCases & FLORES-101 is intended for evaluation purposes in machine translation, particularly for low-resource and multilingual machine translation. The dataset enables evaluation of many-to-many multilingual translation systems, as all translations are multilingually aligned. Beyond translation, FLORES-101 can be used to evaluate tasks such as sentence and document classification, language identification, and multilingual domain adaptation. The dataset supports document-level translation evaluation and multimodal translation evaluation via images. The dataset is divided into three splits: dev (for hyper-parameter tuning), devtest (for testing during development phase), and test (available via evaluation server). \\
\hline
\end{longtable}

\newpage

\section*{3. CIFAR (CIFAR-10/100)}
\textbf{Fill Rate:} 43.3\% (13/30 fields) \quad \textbf{Time:} 23.47s

\begin{longtable}{|p{4cm}|p{12cm}|}
\hline
\textbf{Field} & \textbf{Value} \\
\hline
\endfirsthead
\hline
\textbf{Field} & \textbf{Value} \\
\hline
\endhead
\hline
\endfoot

sc:name & CIFAR-10 Fast Training Methods (airbench) \\
\hline
sc:description & Training methods for CIFAR-10 which reach 94\% accuracy in 3.29 seconds, 95\% in 10.4 seconds, and 96\% in 46.3 seconds on a single NVIDIA A100 GPU. The methods include airbench94\_compiled.py, airbench94.py, airbench95.py, and airbench96.py. \\
\hline
sc:url & https://github.com/KellerJordan/cifar10-airbench \\
\hline
sc:license & null \\
\hline
sc:creator & Keller Jordan \\
\hline
sc:publisher & null \\
\hline
sc:datePublished & 2024-04-05 \\
\hline
sc:inLanguage & en \\
\hline
cr:citeAs & arXiv:2404.00498v2 [cs.LG] 5 Apr 2024 \\
\hline
cr:isLiveDataset & null \\
\hline
rai:dataCollection & The CIFAR-10 dataset contains 60,000 32x32 color images, each labeled as one of ten classes. It is divided into a training set of 50,000 images and a validation set of 10,000 images. \\
\hline
rai:dataCollectionType & Secondary Data analysis \\
\hline
rai:dataCollectionMissingData & null \\
\hline
rai:dataCollectionRawData & CIFAR-10 dataset (Krizhevsky et al., 2009) \\
\hline
rai:dataCollectionTimeframe & null \\
\hline
rai:dataImputationProtocol & null \\
\hline
rai:dataManipulationProtocol & null \\
\hline
rai:dataPreprocessingProtocol & Images are normalized using CIFAR\_MEAN = (0.4914, 0.4822, 0.4465) and CIFAR\_STD = (0.2470, 0.2435, 0.2616). Images are converted from uint8 to half precision (fp16) and divided by 255. Images are permuted to channels\_last memory format. For augmentation, random horizontal flipping and 2-pixel random translation with reflection padding are used. A frozen patch-whitening initialization is applied to the first convolutional layer using eigenvectors of the covariance matrix of 2x2 patches across the training distribution. \\
\hline
rai:dataAnnotationProtocol & null \\
\hline
rai:dataAnnotationPlatform & null \\
\hline
rai:dataAnnotationAnalysis & null \\
\hline
rai:annotationsPerItem & null \\
\hline
rai:annotatorDemographics & null \\
\hline
rai:machineAnnotationTools & null \\
\hline
rai:dataReleaseMaintenancePlan & null \\
\hline
rai:personalSensitiveInformation & null \\
\hline
rai:dataSocialImpact & null \\
\hline
rai:dataBiases & null \\
\hline
rai:dataLimitations & The methods are specifically optimized for CIFAR-10 and may require hyperparameter tuning for other datasets. The training methods use test-time augmentation which improves performance but increases inference time. Without TTA, the three training methods attain 93.2\%, 94.4\%, and 95.6\% mean accuracy respectively. \\
\hline
rai:dataUseCases & Training: The methods are designed for training neural networks on CIFAR-10 to accelerate research and reduce experimental costs. The methods can be used for projects involving massive numbers of trained networks, such as data attribution studies and training variance studies. Usage Guidelines: The methods are intended for experiments where many networks are trained at once to amortize compilation costs (for compiled version). The non-compiled variant can be easily installed and run. The methods support scenarios requiring different performance levels (94\%, 95\%, 96\% accuracy). \\
\hline
\end{longtable}

\newpage

\section*{4. Visual Genome}
\textbf{Fill Rate:} 70.0\% (21/30 fields) \quad \textbf{Time:} 43.43s

\begin{longtable}{|p{4cm}|p{12cm}|}
\hline
\textbf{Field} & \textbf{Value} \\
\hline
\endfirsthead
\hline
\textbf{Field} & \textbf{Value} \\
\hline
\endhead
\hline
\endfoot

sc:name & Visual Genome \\
\hline
sc:description & Visual Genome is a dataset to enable the modeling of relationships between objects in images. It contains over 100K images where each image has an average of 21 objects, 18 attributes, and 18 pairwise relationships between objects. The dataset includes dense annotations of objects, attributes, and relationships within each image, region descriptions, and question answer pairs. \\
\hline
sc:url & https://visualgenome.org \\
\hline
sc:license & null \\
\hline
sc:creator & Ranjay Krishna, Yuke Zhu, Oliver Groth, Justin Johnson, Kenji Hata, Joshua Kravitz, Stephanie Chen, Yannis Kalantidis, Li-Jia Li, David A. Shamma, Michael S. Bernstein, Li Fei-Fei \\
\hline
sc:publisher & null \\
\hline
sc:datePublished & 2016-02-23 \\
\hline
sc:inLanguage & en \\
\hline
cr:citeAs & Ranjay Krishna et al. Visual Genome: Connecting Language and Vision Using Crowdsourced Dense Image Annotations. arXiv:1602.07332v1 [cs.CV] 23 Feb 2016 \\
\hline
cr:isLiveDataset & null \\
\hline
rai:dataCollection & Visual Genome uses 108,249 images from the intersection of the YFCC100M and MS-COCO datasets. The images are real-world, non-iconic images that were uploaded onto Flickr by users. Data collection involved crowdsourcing through Amazon Mechanical Turk where workers drew bounding boxes and wrote descriptions for regions, extracted objects, attributes, and relationships, and created question-answer pairs. \\
\hline
rai:dataCollectionType & Crowdsourcing, Manual Human Curator, Web Scraping, Secondary Data analysis \\
\hline
rai:dataCollectionMissingData & null \\
\hline
rai:dataCollectionRawData & The dataset uses 108,249 images from the intersection of the YFCC100M (Thomee et al., 2016) and MS-COCO (Lin et al., 2014) datasets. These images are real-world, non-iconic images that were uploaded onto Flickr by users. \\
\hline
rai:dataCollectionTimeframe & The dataset was collected over the course of 6 months after 15 months of experimentation and iteration on the data representation. \\
\hline
rai:dataImputationProtocol & null \\
\hline
rai:dataManipulationProtocol & null \\
\hline
rai:dataPreprocessingProtocol & For region descriptions: lowercase, lemmatize, and strip excess whitespace from all attributes and relationships. For objects: Stanford NLP tools were used to extract noun phrases from region descriptions and QAs. For attributes: normalize each attribute based on morphology (stemming) and map them to WordNet adjectives with 15 hand-crafted rules. For relationships: ignore all prepositions and try to find WordNet synsets whose sentence frames match with the context of the relationship, with 20 hand-mapped rules. All objects, attributes, relationships, and noun phrases in region descriptions and QAs are canonicalized to WordNet synsets. \\
\hline
rai:dataAnnotationProtocol & Workers were asked to draw three bounding boxes and write three descriptions for regions. The process was repeated until 50 region descriptions were collected per image. Workers were shown the top seven most similar descriptions using BLEU scores to prevent repetition. For objects: workers extracted all objects from descriptions and grounded each object as a bounding box. For attributes and relationships: workers added attributes to objects or connected pairs of objects with relationships based on the text of the description. For QA: workers wrote pairs of questions and answers about images, following rules to start with one of the seven Ws, avoid ambiguous questions, and be precise. Two types of QAs were collected: freeform QAs (8 pairs per image) and region-based QAs (around 20 pairs per image). \\
\hline
rai:dataAnnotationPlatform & Amazon Mechanical Turk \\
\hline
rai:dataAnnotationAnalysis & All Visual Genome data go through a verification stage using two strategies: majority voting (three unique workers vote on whether annotation is factually correct, requiring at least two out of three to verify) and rapid judgments (used only for objects verification). Annotations are added to the dataset if verified as correct. \\
\hline
rai:annotationsPerItem & Each image has an average of 42 region descriptions, 21 objects, 16 attributes, 18 relationships, and 17 question-answer pairs. \\
\hline
rai:annotatorDemographics & A total of over 33,000 unique workers contributed to the dataset. 93.02\% of workers contributed from the United States, 1.29\% from Philippines, 1.13\% from Kenya, 0.94\% from India, 0.50\% from Russia, 0.47\% from Canada, and 2.65\% from other countries. The majority of workers were between the ages of 25 and 34 years old, with the youngest contributor being 18 years old and the oldest being 68 years old. The dataset had a near-balanced split of 54.15\% male and 45.85\% female workers. \\
\hline
rai:machineAnnotationTools & Stanford NLP tools (Manning et al., 2014) were used to extract nouns automatically from region descriptions. Stanford's dependency parser was used to extract nouns and send them to workers as suggestions. VGG 16-layer network (Simonyan and Zisserman, 2014) was used for feature extraction in experiments. \\
\hline
rai:dataReleaseMaintenancePlan & null \\
\hline
rai:personalSensitiveInformation & null \\
\hline
rai:dataSocialImpact & null \\
\hline
rai:dataBiases & The dataset is somewhat biased towards images of people, as the top synsets refer to images of people playing sports. The images contain a bias towards certain sports (skiing, surfboarding) due to the source images. Common visual elements like 'green grass,' 'tree [in] distance,' and 'blue sky' occur much more often than other, more nuanced elements. \\
\hline
rai:dataLimitations & The canonicalization mappings are not perfect and still contain some ambiguity. The model for attribute prediction tends to associate objects with attributes (e.g. elephant with grazing). For relationship prediction, there is high intra-class variability of most relationships, making them not visually distinct enough for discriminative models to learn effectively. Current models fail to describe all the different events captured in an image and instead provide only a high-level summary. \\
\hline
rai:dataUseCases & The dataset can be used for training and benchmarking models for comprehensive scene understanding tasks including: attribute prediction, relationship prediction, dense image captioning, visual question answering, image understanding evaluation, relationship extraction, semantic image retrieval, object detection, and scene graph generation. The dataset enables research on cognitive computer vision tasks that require understanding interactions and relationships between objects in images. \\
\hline
\end{longtable}

\newpage

\section*{5. MSCOCO (MS COCO)}
\textbf{Fill Rate:} 66.7\% (20/30 fields) \quad \textbf{Time:} 48.95s

\begin{longtable}{|p{4cm}|p{12cm}|}
\hline
\textbf{Field} & \textbf{Value} \\
\hline
\endfirsthead
\hline
\textbf{Field} & \textbf{Value} \\
\hline
\endhead
\hline
\endfoot

sc:name & Microsoft COCO: Common Objects in Context \\
\hline
sc:description & A large-scale dataset for object detection and segmentation containing photos of 91 object types in complex everyday scenes. The dataset contains 2.5 million labeled instances in 328k images with per-instance segmentations to aid in precise object localization. \\
\hline
sc:url & http://mscoco.org/ \\
\hline
sc:license & null \\
\hline
sc:creator & Tsung-Yi Lin, Michael Maire, Serge Belongie, Lubomir Bourdev, Ross Girshick, James Hays, Pietro Perona, Deva Ramanan, C. Lawrence Zitnick, Piotr Doll\'{a}r \\
\hline
sc:publisher & null \\
\hline
sc:datePublished & 2015-02-21 \\
\hline
sc:inLanguage & en \\
\hline
cr:citeAs & T. Lin, M. Maire, S. Belongie, J. Hays, P. Perona, D. Ramanan, P. Doll\'{a}r, and C. L. Zitnick, "Microsoft COCO: Common objects in context," in ECCV, 2014. \\
\hline
cr:isLiveDataset & null \\
\hline
rai:dataCollection & Images were collected from Flickr using a novel pipeline that queries for pairs of objects in conjunction with scene-based queries to gather non-iconic images containing contextual relationships. The dataset was created by searching for pairwise combinations of object categories (e.g., 'dog + car') and scene/object category pairs rather than single object categories to avoid iconic images. At most 5 photos taken by a single photographer within a short time window were downloaded. For categories where enough images could not be found, single category searches were performed with explicit filtering to remove iconic images. \\
\hline
rai:dataCollectionType & Web Scraping, Manual Human Curator \\
\hline
rai:dataCollectionMissingData & null \\
\hline
rai:dataCollectionRawData & Images were harvested from Flickr, which contains photos uploaded by amateur photographers with searchable metadata and keywords. Images were retrieved using pairwise object category queries, scene/object category pairs, and in rare cases single object category queries with filtering. \\
\hline
rai:dataCollectionTimeframe & null \\
\hline
rai:dataImputationProtocol & null \\
\hline
rai:dataManipulationProtocol & Near-duplicate images were explicitly removed using detection methods. Images were grouped by photographer and date taken to minimize the chance of near-duplicate images existing across train/validation/test splits. \\
\hline
rai:dataPreprocessingProtocol & An explicit filtering stage was performed to remove iconic images when single category searches were used. Workers were shown grids of 128 candidate images and asked to remove invalid or iconic images. Instructions with examples of iconic and non-iconic images were provided to workers. \\
\hline
rai:dataAnnotationProtocol & The annotation pipeline consisted of three primary stages: (1) Category labeling: Workers used a hierarchical approach to indicate which of 91 object categories were present in each image by dragging category icons onto instances. Categories were grouped into 11 super-categories. 8 workers labeled each image, and a category was considered present if any worker indicated it. (2) Instance spotting: Workers placed crosses on all instances (up to 10) of each labeled category. The location from the previous stage was shown to prime workers. A magnifying glass feature was provided for small objects. 8 workers labeled each image. (3) Instance segmentation: Workers segmented each object instance using a modified interface from the OpenSurfaces project. Workers completed a training task for each object category before being allowed to segment instances. Segmentations were verified by 3-5 workers, requiring at least 4 of 5 favorable votes to be accepted. For images with more than 10 instances of a category, remaining instances were marked as 'crowds' using a single segment. \\
\hline
rai:dataAnnotationPlatform & Amazon Mechanical Turk (AMT) \\
\hline
rai:dataAnnotationAnalysis & Category presence was determined by union of 8 annotators to maximize recall. False positives were handled in subsequent stages. For segmentation verification, each segmentation was initially shown to 3 annotators. If any indicated it was bad, it was shown to 2 additional workers. Segmentations not receiving at least 4 of 5 favorable votes were discarded. Worker quality was analyzed by comparing precision and recall of AMT workers with expert workers (co-authors). The union of 8 AMT workers achieved greater recall than any individual expert worker. Over 99\% of all object categories not later rejected as false positives were estimated to be detected given 8 annotators. \\
\hline
rai:annotationsPerItem & Category labeling: 8 workers per image. Instance spotting: 8 workers per image. Instance segmentation: 1 worker per instance. Segmentation verification: 3-5 workers per segmentation (requiring 4 of 5 favorable votes). \\
\hline
rai:annotatorDemographics & null \\
\hline
rai:machineAnnotationTools & null \\
\hline
rai:dataReleaseMaintenancePlan & The MS COCO dataset is split into two roughly equal parts. The first half (2014 release) contains 82,783 training, 40,504 validation, and 40,775 testing images. The cumulative 2015 release will contain a total of 165,482 train, 81,208 val, and 81,434 test images. Annotations for train and validation data are released, but not for test. The 2014 release is limited to a subset of 80 categories (11 categories omitted). Some of these categories may be added in the cumulative 2015 release. The dataset will evolve and grow over time with up-to-date information available online. \\
\hline
rai:personalSensitiveInformation & null \\
\hline
rai:dataSocialImpact & null \\
\hline
rai:dataBiases & The dataset aims to reduce bias toward iconic images by collecting non-iconic images containing objects in natural contexts with non-canonical viewpoints. However, the selection of 91 entry-level object categories may reflect certain cultural or practical biases. Categories were selected based on votes from co-authors considering how commonly they occur, their usefulness for practical applications, and their diversity relative to other categories. \\
\hline
rai:dataLimitations & The dataset is limited to 'thing' categories (objects with clear boundaries) and does not include 'stuff' categories (materials and objects with no clear boundaries like sky, street, grass), though such categories could provide significant contextual information. The dataset uses entry-level categories rather than fine-grained categories to enable collection of significant instances per category. Object-part categories (face, hands, wheels) are not included but could be beneficial. For images with more than 10-15 instances of a category, remaining instances are marked as 'crowds' and ignored in evaluation. Some categories (hat, shoe, eyeglasses, mirror, window, door, street sign, plate, desk, blender, hair brush) were omitted from the 2014 release due to various issues including too many instances, ambiguity, confusion with other categories, or too few instances. \\
\hline
rai:dataUseCases & The dataset is designed for three core research problems in scene understanding: detecting non-iconic views of objects, contextual reasoning between objects, and precise 2D localization of objects. Primary use cases include object detection and instance segmentation. The dataset can be used for training, validation, and testing of object recognition algorithms. Baseline performance analysis is provided for bounding box and segmentation detection using Deformable Parts Model. The dataset enables evaluation of both detection accuracy and segmentation quality. Areas marked as 'crowds' are ignored and do not affect a detector's score during evaluation. \\
\hline
\end{longtable}

\newpage

\section*{6. MMLU}
\textbf{Fill Rate:} 66.7\% (20/30 fields) \quad \textbf{Time:} 33.54s

\begin{longtable}{|p{4cm}|p{12cm}|}
\hline
\textbf{Field} & \textbf{Value} \\
\hline
\endfirsthead
\hline
\textbf{Field} & \textbf{Value} \\
\hline
\endhead
\hline
\endfoot

sc:name & Measuring Massive Multitask Language Understanding (MMLU) \\
\hline
sc:description & A test to measure a text model's multitask accuracy covering 57 tasks including elementary mathematics, US history, computer science, law, and more. The test spans subjects in STEM, the humanities, the social sciences, and other areas at difficulty levels from elementary to advanced professional. To attain high accuracy on this test, models must possess extensive world knowledge and problem solving ability. \\
\hline
sc:url & github.com/hendrycks/test \\
\hline
sc:license & null \\
\hline
sc:creator & Dan Hendrycks, Collin Burns, Steven Basart, Andy Zou, Mantas Mazeika, Dawn Song, Jacob Steinhardt \\
\hline
sc:publisher & ICLR 2021 \\
\hline
sc:datePublished & 2021 \\
\hline
sc:inLanguage & en \\
\hline
cr:citeAs & Hendrycks et al., 2021, Measuring Massive Multitask Language Understanding, ICLR 2021 \\
\hline
cr:isLiveDataset & false \\
\hline
rai:dataCollection & Questions were manually collected by graduate and undergraduate students from freely available sources online. These include practice questions for tests such as the Graduate Record Examination and the United States Medical Licensing Examination. It also includes questions designed for undergraduate courses and questions designed for readers of Oxford University Press books. \\
\hline
rai:dataCollectionType & Manual Human Curator \\
\hline
rai:dataCollectionMissingData & null \\
\hline
rai:dataCollectionRawData & Practice questions for tests such as the Graduate Record Examination and the United States Medical Licensing Examination, questions designed for undergraduate courses, and questions designed for readers of Oxford University Press books. Sources include freely available questions online from PDFs or websites where questions and answers are on separate pages. \\
\hline
rai:dataCollectionTimeframe & null \\
\hline
rai:dataImputationProtocol & null \\
\hline
rai:dataManipulationProtocol & null \\
\hline
rai:dataPreprocessingProtocol & Questions were collected and formatted as multiple-choice questions with four options (A, B, C, D). The dataset was split into a few-shot development set (5 questions per subject), a validation set (1540 questions), and a test set (14079 questions). Each subject contains 100 test examples at the minimum. Mathematics expressions were encoded using LaTeX or symbols such as * and \^{} for multiplication and exponentiation respectively. \\
\hline
rai:dataAnnotationProtocol & Questions were manually collected by graduate and undergraduate students from freely available sources. Ground truth answers were obtained from the original source materials. \\
\hline
rai:dataAnnotationPlatform & null \\
\hline
rai:dataAnnotationAnalysis & null \\
\hline
rai:annotationsPerItem & null \\
\hline
rai:annotatorDemographics & Graduate and undergraduate students \\
\hline
rai:machineAnnotationTools & null \\
\hline
rai:dataReleaseMaintenancePlan & The test and code is available at github.com/hendrycks/test. To reduce the probability that future models encounter exact questions during test-time, a list of question sources will be provided. \\
\hline
rai:personalSensitiveInformation & null \\
\hline
rai:dataSocialImpact & The test can help researchers pinpoint important shortcomings of models, making it easier to gain a clearer picture of state-of-the-art capabilities. Models perform especially poorly on socially relevant subjects including morality and law. For future systems to be aligned with human values, high performance on these tasks is crucial. \\
\hline
rai:dataBiases & Models have lopsided performance with GPT-3 having almost 70\% accuracy for its best subject but near-random performance for several other subjects. Models are notably poor at modeling human (dis)approval, as evident by the low performance on the Professional Law and Moral Scenarios tasks. Models also have difficulty performing calculations, exhibiting poor performance on Elementary Mathematics and many other STEM subjects with calculation-heavy problems. \\
\hline
rai:dataLimitations & State-of-the-art models still struggle at learning and applying knowledge from pretraining. Models do not match expert-level performance (90\%) on any subject. Models are uncalibrated with confidence being up to 24\% off from actual accuracy. Models have difficulty with tasks that require calculations. Current models perform especially poorly on socially relevant subjects including morality and law. GPT-3 acquires knowledge quite unlike humans, learning about topics in a pedagogically unusual order (e.g., does better on College Medicine than Elementary Mathematics). The test is text-only and does not incorporate multimodal information. \\
\hline
rai:dataUseCases & The benchmark is designed to measure knowledge acquired during pretraining by evaluating models exclusively in zero-shot and few-shot settings. It can be used to analyze models across many tasks and to identify important shortcomings. The test assesses language understanding in greater breadth and depth than previous benchmarks. Models are assessed in zero-shot, few-shot, or transfer settings. The dev set is used for few-shot prompts, the val set could be used for hyperparameter tuning, and the test set is used to compute the final accuracy. \\
\hline
\end{longtable}

\newpage

\section*{7. MMMU}
\textbf{Fill Rate:} 56.7\% (17/30 fields) \quad \textbf{Time:} 38.90s

\begin{longtable}{|p{4cm}|p{12cm}|}
\hline
\textbf{Field} & \textbf{Value} \\
\hline
\endfirsthead
\hline
\textbf{Field} & \textbf{Value} \\
\hline
\endhead
\hline
\endfoot

sc:name & MMMU: A Massive Multi-discipline Multimodal Understanding and Reasoning Benchmark for Expert AGI \\
\hline
sc:description & MMMU is a benchmark designed to evaluate multimodal models on massive multi-discipline tasks demanding college-level subject knowledge and deliberate reasoning. It includes 11.5K meticulously collected multimodal questions from college exams, quizzes, and textbooks, covering six core disciplines: Art \& Design, Business, Science, Health \& Medicine, Humanities \& Social Science, and Tech \& Engineering. These questions span 30 subjects and 183 subfields, comprising 30 highly heterogeneous image types. \\
\hline
sc:url & https://mmmu-benchmark.github.io/ \\
\hline
sc:license & null \\
\hline
sc:creator & Xiang Yue, Yuansheng Ni, Kai Zhang, Tianyu Zheng, Ruoqi Liu, Ge Zhang, Samuel Stevens, Dongfu Jiang, Weiming Ren, Yuxuan Sun, Cong Wei, Botao Yu, Ruibin Yuan, Renliang Sun, Ming Yin, Boyuan Zheng, Zhenzhu Yang, Yibo Liu, Wenhao Huang, Huan Sun, Yu Su, Wenhu Chen \\
\hline
sc:publisher & null \\
\hline
sc:datePublished & 2023-11-27 \\
\hline
sc:inLanguage & en \\
\hline
cr:citeAs & null \\
\hline
cr:isLiveDataset & null \\
\hline
rai:dataCollection & Data is primarily collected from free online resources, quizzes, textbooks, and other study materials. The collection takes three stages: (1) selecting 30 subjects from six different disciplines based on the principle that visual inputs should be commonly adopted in the subjects; (2) recruiting over 50 university students, including co-authors, specializing in these majors as annotators to assist in question collection from major textbooks and online resources, creating new questions based on their expertise where necessary; (3) data quality control involving duplicate identification, format and typo checking, and difficulty categorization. \\
\hline
rai:dataCollectionType & Manual Human Curator, Secondary Data analysis, Document analysis \\
\hline
rai:dataCollectionMissingData & null \\
\hline
rai:dataCollectionRawData & Questions were manually collected by a team of 50 college students (including coauthors) from various disciplines and subjects, drawing from online sources, textbooks, and lecture materials. Sources include free online resources, quizzes, major textbooks, and other study materials. \\
\hline
rai:dataCollectionTimeframe & null \\
\hline
rai:dataImputationProtocol & null \\
\hline
rai:dataManipulationProtocol & null \\
\hline
rai:dataPreprocessingProtocol & Data cleaning was performed in two steps: (1) lexical overlap and source URL similarity were employed to identify potential duplicate problems, which were then reviewed by the authors to identify and eliminate any duplications; (2) problems were distributed among different co-authors for format and typo checking to ensure adherence to a standardized format, undertaking necessary corrections where deviations were found; (3) authors categorized the problems into four difficulty levels (very easy, easy, medium, and hard), with approximately 10\% of the problems classified as very easy being excluded from the benchmark. \\
\hline
rai:dataAnnotationProtocol & Annotators were instructed to collect multimodal questions from major textbooks and online resources, creating new questions based on their expertise where necessary. All questions must contain one or more images, be written in English, meet college-level difficulty, not be ambiguous, and be answerable with one of the given options or a short answer. Annotators categorized each question as either multiple-choice or open-ended and annotated all fields including the question, answer options for multiple-choice questions, the correct answer, image types, question difficulty, and explanation (if available). \\
\hline
rai:dataAnnotationPlatform & A user-friendly web interface was developed for data annotation and verification. \\
\hline
rai:dataAnnotationAnalysis & null \\
\hline
rai:annotationsPerItem & null \\
\hline
rai:annotatorDemographics & Over 50 university students, including co-authors, specializing in various majors across the 30 subjects covered in the benchmark. Additionally, 90 college senior students were selected to represent a wide range of experts in the corresponding 30 subjects (3 student experts per subject) for human expert evaluation. \\
\hline
rai:machineAnnotationTools & null \\
\hline
rai:dataReleaseMaintenancePlan & null \\
\hline
rai:personalSensitiveInformation & Annotators were instructed to avoid collecting questions that contain any private information. \\
\hline
rai:dataSocialImpact & The benchmark is designed to push the boundaries of what large multimodal models can achieve and is intended to be instrumental in developing next-generation multimodal foundation models and monitoring the progress towards Expert AGI. The authors note that strong performance on MMMU should be a necessary criterion for an Expert AGI system, which marks a critical milestone denoting an AI system that reaches at least 90th percentile of skilled adults in a broad range of tasks, potentially leading to significant risks of job displacement and economic disruption. \\
\hline
rai:dataBiases & null \\
\hline
rai:dataLimitations & The manual curation process may carry biases, and the focus on college-level subjects might not be sufficient for testing Expert AGI. The benchmark combines multiple-choice questions with concise open-ended questions to strike a balance between complexity and practicality, enabling the assessment of diverse subjects while addressing the challenges associated with evaluating open-ended responses. MMMU is not a sufficient test for Expert AGI as there lacks a direct mapping between performance on MMMU and 90th percentile of skilled adults, nor are college exams the only tasks an AGI shall tackle. \\
\hline
rai:dataUseCases & The benchmark is designed for evaluating multimodal models on massive multi-discipline tasks demanding college-level subject knowledge and deliberate reasoning. It is intended for testing, validation, and benchmarking purposes to assess the capabilities of large multimodal models in expert-level perception and reasoning with domain-specific knowledge. The dataset is divided into a few-shot development set (5 questions per subject), a validation set (approximately 900 questions for hyperparameter selection), and a test set (10.5K questions). \\
\hline
\end{longtable}

\newpage

\section*{8. MathVista}
\textbf{Fill Rate:} 76.7\% (23/30 fields) \quad \textbf{Time:} 53.69s

\begin{longtable}{|p{4cm}|p{12cm}|}
\hline
\textbf{Field} & \textbf{Value} \\
\hline
\endfirsthead
\hline
\textbf{Field} & \textbf{Value} \\
\hline
\endhead
\hline
\endfoot

sc:name & MATHVISTA \\
\hline
sc:description & MATHVISTA is a benchmark designed to combine challenges from diverse mathematical and visual tasks. It consists of 6,141 examples, derived from 28 existing multimodal datasets involving mathematics and 3 newly created datasets (i.e., IQTest, FunctionQA, and PaperQA). Completing these tasks requires fine-grained, deep visual understanding and compositional reasoning. \\
\hline
sc:url & https://mathvista.github.io \\
\hline
sc:license & null \\
\hline
sc:creator & Pan Lu, Hritik Bansal, Tony Xia, Jiacheng Liu, Chunyuan Li, Hannaneh Hajishirzi, Hao Cheng, Kai-Wei Chang, Michel Galley, Jianfeng Gao \\
\hline
sc:publisher & UCLA, University of Washington, Microsoft Research, Redmond \\
\hline
sc:datePublished & 2024 \\
\hline
sc:inLanguage & English (93.43\%), Chinese, Persian \\
\hline
cr:citeAs & Pan Lu, Hritik Bansal, Tony Xia, Jiacheng Liu, Chunyuan Li, Hannaneh Hajishirzi, Hao Cheng, Kai-Wei Chang, Michel Galley, Jianfeng Gao. MATHVISTA: Evaluating Mathematical Reasoning of Foundation Models in Visual Contexts. ICLR 2024. \\
\hline
cr:isLiveDataset & null \\
\hline
rai:dataCollection & MATHVISTA incorporates 28 existing multimodal datasets, including 9 math-targeted question answering (MathQA) datasets and 19 VQA datasets. Three new datasets (IQTest, FunctionQA, PaperQA) were created. For MathQA datasets, each source dataset is limited to up to 400 examples to ensure balanced representation. For VQA datasets, heuristic rules were used to automatically select examples likely to involve mathematical reasoning from a large pool of candidates, followed by manual labeling by three expert annotators to determine if they involve mathematical reasoning. \\
\hline
rai:dataCollectionType & Manual Human Curator, Secondary Data analysis \\
\hline
rai:dataCollectionMissingData & null \\
\hline
rai:dataCollectionRawData & Data is sourced from 28 existing multimodal datasets (9 MathQA datasets and 19 VQA datasets) and 3 newly created datasets. MathQA datasets include GPS, MWP, and TQA datasets. VQA datasets were reviewed from more than 70 datasets, collecting 19 that contain math-related instances. New datasets include IQTest (228 examples from puzzle test figures on online learning platforms), FunctionQA (400 examples emphasizing functional plots), and PaperQA (107 examples from academic papers released in August 2023 on Huggingface). \\
\hline
rai:dataCollectionTimeframe & New datasets (PaperQA) sourced from papers released in August 2023. Evaluation dates for Multimodal Bard range from Sep 8, 2023 to Sep 10, 2023. Evaluation dates for GPT-4V range from Oct 7, 2023 to Oct 15, 2023. \\
\hline
rai:dataImputationProtocol & null \\
\hline
rai:dataManipulationProtocol & null \\
\hline
rai:dataPreprocessingProtocol & For VQA datasets, heuristic rules were designed to automatically select examples likely to involve mathematical reasoning. Examples with numeric answers or those containing quantity words were selected. This automatic filtration yielded 4,949 VQA-format examples. Each source dataset is limited to up to 400 examples to ensure balanced representation. Questions are normalized to required answer formats (e.g., an option letter or an integer). \\
\hline
rai:dataAnnotationProtocol & All questions for the three new datasets were manually annotated by graduate students in STEM fields and further refined through a rigorous review process. Each dataset was independently annotated by three reviewers, resulting in a high inter-annotation consistency rate of 99.2\%. Among the newly collected 736 questions, only 6 exhibited disagreements in the annotated answers. Discrepancies were resolved through discussion among the entire review team. For VQA datasets, three expert annotators manually labeled examples to determine if they involve mathematical reasoning, utilizing majority voting. Metadata annotations (question type, answer type, language, source, category, task, grade level, visual context, mathematical reasoning types) were obtained from source dataset details. Expert annotators manually labeled mathematical reasoning categories for 1,000 examples to verify automatic annotation quality, showing 94.1\% exact match and 98.79\% individual label accuracy. \\
\hline
rai:dataAnnotationPlatform & Custom annotation tool with GUI (illustrated in Figure 23 and Figure 24 in the paper) \\
\hline
rai:dataAnnotationAnalysis & A two-step annotation process was employed. Initially, each dataset was independently annotated by three reviewers, resulting in a high inter-annotation consistency rate of 99.2\%. Specifically, among the newly collected 736 questions, only 6 exhibited disagreements in the annotated answers. Discrepancies were resolved through discussion among the entire review team. For mathematical reasoning labeling, Fleiss Kappa score of 0.775 was achieved, indicating substantial consistency among annotators. For VQA dataset selection, majority voting was utilized among three expert annotators. \\
\hline
rai:annotationsPerItem & For the three new datasets, each example was independently annotated by three reviewers. For VQA dataset selection, three expert annotators labeled each example. For human performance evaluation, each question from the testmini subset was assigned to five annotators. \\
\hline
rai:annotatorDemographics & Annotators for the three new datasets were graduate students in STEM fields. For human performance evaluation on Amazon Mechanical Turk, eligible annotators must have a satisfactory annotating history, successfully pass qualification examples, and possess a high school degree or higher. Results were retained exclusively from annotators who had achieved a high school diploma or higher. Annotators have a history of completing more than 5,000 HIT tasks and boast an acceptance score higher than 0.99. \\
\hline
rai:machineAnnotationTools & GPT-4 was used as an answer extractor for evaluation, with more than 99.5\% accuracy in extracting answer text from detailed responses. EasyOCR was used for OCR text detection. Multimodal Bard was used to generate image captions for augmented LLMs. \\
\hline
rai:dataReleaseMaintenancePlan & MATHVISTA consists of 6,141 examples, divided into two subsets: testmini (1,000 examples) and test (5,141 examples). The answer labels for test will not be publicly released to prevent data contamination, and an online evaluation platform will be maintained. The authors are committed to maintaining the datasets, such as refining potential data noise, in response to community feedback and evolving the leaderboard in response to new models. \\
\hline
rai:personalSensitiveInformation & null \\
\hline
rai:dataSocialImpact & MATHVISTA is designed to facilitate the development of general-purpose AI agents capable of tackling mathematically intensive and visually rich real-world tasks. It has wide-ranging applications including solving complex problems in educational disciplines, helping analysts with logical queries about statistical data, and assisting in theorem proving and scientific discovery in advanced research fields. \\
\hline
rai:dataBiases & null \\
\hline
rai:dataLimitations & While MATHVISTA encompasses a broad spectrum of tasks and visual contexts, there may be gaps in the representation of certain types of mathematical problems and visuals. The dataset's focus on mathematical reasoning within visual contexts, spanning specific domains like science and college-level math, necessitates a more labor-intensive process for collecting high-quality data. Only 85.6\% of examples have annotations due to sourcing from original data providers. Annotations lack a unified format and structure due to heterogeneity of sources. The scalability and generalizability of the benchmark to other domains remain a concern. \\
\hline
rai:dataUseCases & MATHVISTA is intended for evaluating mathematical reasoning capabilities of foundation models in visual contexts. It can be used for Testing and Validation of Large Language Models (LLMs) and Large Multimodal Models (LMMs). The testmini subset (1,000 examples) is intended for model development validation or for those with limited computing resources. The benchmark is designed to assess models across multiple tasks (figure question answering, geometry problem solving, math word problem, textbook question answering, visual question answering) and mathematical reasoning types (algebraic, arithmetic, geometry, logical, numeric common sense, scientific, statistical reasoning). \\
\hline
\end{longtable}

\end{document}
